\documentclass[a4paper,11pt]{article}
\usepackage[utf8]{inputenc}
\usepackage[T1]{fontenc}
\usepackage[french]{babel}
\usepackage[left=2.5cm,right=2.5cm,top=1.8cm,bottom=1.8cm]{geometry}
\usepackage{hyperref}
\usepackage{xcolor}
\usepackage{fontawesome5}
\usepackage{lmodern}

% --- Couleurs ---
\definecolor{primary}{RGB}{0, 82, 147}

\hypersetup{colorlinks=true, linkcolor=primary, urlcolor=primary}

\begin{document}

\noindent
\textbf{Loïc Aron MBASSI EWOLO} \\
Élève-Ingénieur ENSEA / Polytechnique Yaoundé \\
Cergy, France \\
\faPhone\ (+33) 7 59 52 33 98 \quad \faEnvelope\ \href{mailto:loicaron.mbassiewolo@ensea.fr}{loicaron.mbassiewolo@ensea.fr}

\vspace{0.6cm}

\noindent
\textbf{CEA Gramat} \\
Service de Recrutement \\
Gramat (46)

\hfill Cergy, le 23 janvier 2026

\vspace{0.5cm}

\noindent\textbf{Objet : Stage Développement d'outils de génération de microstructures}

\vspace{0.4cm}

Madame, Monsieur,

\vspace{0.3cm}

Je suis en deuxième année à l'ENSEA dans le cadre d'un double diplôme avec l'École Polytechnique de Yaoundé, et je cherche un stage de 5 à 6 mois. Votre offre m'a interpellé parce qu'elle touche à un sujet qui me plaît : développer des algorithmes qui génèrent quelque chose de concret, ici des microstructures.

\vspace{0.25cm}

Ce qui m'a attiré dans cette offre, c'est le problème du placement de grains avec contrainte de fraction volumique. Ça me rappelle un projet sur lequel j'ai travaillé : l'optimisation de routes de collecte de déchets. On devait placer des points de passage tout en respectant des contraintes de capacité et de distance. On a utilisé une approche inspirée du voyageur de commerce, et j'ai passé pas mal de temps à optimiser le code pour que ça tourne en temps raisonnable. Le projet a gagné le premier prix au hackathon CITS, mais surtout j'ai appris qu'en optimisation, les détails comptent énormément.

\vspace{0.25cm}

Sur la partie IA pour augmenter la base de données de grains, j'ai une expérience qui peut être utile. Pour un projet de reconnaissance de langue des signes, j'ai dû créer mon propre dataset de gestes. Comme on n'avait pas assez de données, j'ai travaillé sur l'augmentation : rotations, translations, variations d'éclairage. C'est basique, mais ça m'a appris à réfléchir à ce qui fait qu'une donnée synthétique reste réaliste.

\vspace{0.25cm}

Côté technique, je code principalement en Python et en C. J'ai aussi développé un outil open source qui génère du code à partir de diagrammes UML --- ça m'a appris à parser des fichiers XML et à structurer un générateur. Ce n'est pas de la génération de microstructures, mais la logique de « prendre une entrée, appliquer des règles, produire une sortie » est similaire.

\vspace{0.25cm}

Je n'ai pas de background en physique des matériaux. Ma note de 19.25/20 en physique au bac est loin, et je ne prétends pas maîtriser les enjeux des compositions énergétiques. Par contre, je suis curieux, j'aime comprendre le « pourquoi » derrière ce que je code, et je m'adapte vite à de nouveaux domaines.

\vspace{0.25cm}

Je suis disponible à partir d'avril 2026 et mobile pour Gramat.

\vspace{0.4cm}

Cordialement,

\vspace{0.8cm}

\textbf{Loïc Aron MBASSI EWOLO}

\end{document}
